% ====================================================================
% ФАЙЛ: tasks/02_mkt/mkt_thermodynamics_first_law_part1.tex
% ТЕМА: ПЕРВОЕ НАЧАЛО ТЕРМОДИНАМИКИ. РАБОТА ГАЗА. ВНУТРЕННЯЯ ЭНЕРГИЯ
% ПАЧКА 1: ВАРИАНТЫ 1–5
% ====================================================================

% === ЗАДАЧА 1 ===
% Тег: #МКТ #ТЕРМОДИНАМИКА #КИМ_08 #ВАРИАНТ_01
\begin{taskbasic}[code={\label{task:mkt_therm_v1_n8}}]
\textbf{Задача 1.} На рисунке показано, как меняется давление идеального газа в зависимости от его объёма при переходе из состояния 1 в состояние 2, а затем в состояние 3. В ходе процесса 2–3 газ совершил работу $900$~Дж. Чему равна работа газа в процессе 1–2?

\nopagebreak\smallskip
\begin{center}
\begin{tikzpicture}[scale=1.0, every node/.style={font=\footnotesize}]
    \draw[xstep=1, ystep=1, gray!30, thin] (0,0) grid (4,4);
    \draw[-{Stealth[scale=0.8]}, black, thick] (0,0) -- (4.5,0) node[right] {$V$};
    \draw[-{Stealth[scale=0.8]}, black, thick] (0,0) -- (0,4.5) node[above] {$p$};
    
    \coordinate (1) at (1,3);
    \coordinate (2) at (2,3);
    \coordinate (3) at (3,1);
    
    \draw[ultra thick, brandPrimary, ->, >=stealth] (1) -- (1.6,3);
    \draw[ultra thick, brandPrimary] (1.5,3) -- (2);
    \draw[ultra thick, brandPrimary, ->, >=stealth] (2) -- (2.5,2);
    \draw[ultra thick, brandPrimary] (2.4,2.2) -- (3);
    
    \draw[dashed, gray] (0,1) -- (3,1);
    \draw[dashed, gray] (0,3) -- (1,3);
    \draw[dashed, gray] (1,0) -- (1,3);
    \draw[dashed, gray] (2,0) -- (2,3);
    \draw[dashed, gray] (3,0) -- (3,1);
    
    \fill[black] (1) circle (1.5pt) node[above left] {1};
    \fill[black] (2) circle (1.5pt) node[above right] {2};
    \fill[black] (3) circle (1.5pt) node[above right] {3};
    
    \node[left] at (0,1) {$p_0$};
    \node[left] at (0,3) {$2p_0$};
    \node[below] at (1,0) {$V_0$};
    \node[below] at (2,0) {$2V_0$};
    \node[below] at (3,0) {$3V_0$};
    \node[below left] at (0,0) {0};
\end{tikzpicture}
\end{center}

\kim{8}
\end{taskbasic}
\answer{1200}

% === ЗАДАЧА 2 ===
% Тег: #МКТ #ТЕРМОДИНАМИКА #КИМ_08 #ВАРИАНТ_02
\begin{taskbasic}[code={\label{task:mkt_therm_v2_n8}}]
\textbf{Задача 2.} На рисунке показано, как меняется давление идеального газа в зависимости от его объёма при переходе из состояния 1 в состояние 2, а затем в состояние 3. В ходе процесса 1–2 газ совершил работу $500$~Дж. Чему равна работа газа в процессе 2–3?

\nopagebreak\smallskip
\begin{center}
\begin{tikzpicture}[scale=1.0, every node/.style={font=\footnotesize}]
    \draw[xstep=1, ystep=1, gray!30, thin] (0,0) grid (4,4);
    \draw[-{Stealth[scale=0.8]}, black, thick] (0,0) -- (4.5,0) node[right] {$V$};
    \draw[-{Stealth[scale=0.8]}, black, thick] (0,0) -- (0,4.5) node[above] {$p$};
    
    \coordinate (1) at (1,3);
    \coordinate (2) at (2,3);
    \coordinate (3) at (3,1);
    
    \draw[ultra thick, brandPrimary, ->, >=stealth] (1) -- (1.6,3);
    \draw[ultra thick, brandPrimary] (1.5,3) -- (2);
    \draw[ultra thick, brandPrimary, ->, >=stealth] (2) -- (2.5,2);
    \draw[ultra thick, brandPrimary] (2.4,2.2) -- (3);
    
    \draw[dashed, gray] (0,1) -- (3,1);
    \draw[dashed, gray] (0,3) -- (1,3);
    \draw[dashed, gray] (1,0) -- (1,3);
    \draw[dashed, gray] (2,0) -- (2,3);
    \draw[dashed, gray] (3,0) -- (3,1);
    
    \fill[black] (1) circle (1.5pt) node[above left] {1};
    \fill[black] (2) circle (1.5pt) node[above right] {2};
    \fill[black] (3) circle (1.5pt) node[above right] {3};
    
    \node[left] at (0,1) {$p_0$};
    \node[left] at (0,3) {$2p_0$};
    \node[below] at (1,0) {$V_0$};
    \node[below] at (2,0) {$2V_0$};
    \node[below] at (3,0) {$3V_0$};
    \node[below left] at (0,0) {0};
\end{tikzpicture}
\end{center}

\kim{8}
\end{taskbasic}
\answer{375}

% === ЗАДАЧА 3 ===
% Тег: #МКТ #ТЕРМОДИНАМИКА #КИМ_09 #ВАРИАНТ_03
\begin{taskbasic}[code={\label{task:mkt_therm_v3_n9}}]
\textbf{Задача 3.} Один моль разреженного аргона участвует в циклическом процессе 1–2–3–4–1, показанном на рисунке в переменных $p-V$. Из приведённого ниже списка выберите все верные утверждения, характеризующие работу двигателя.

\nopagebreak\smallskip
\begin{center}
\begin{tikzpicture}[scale=1.1, every node/.style={font=\footnotesize}]
    \draw[xstep=1, ystep=1, gray!30, thin] (0,0) grid (4,3);
    \draw[-{Stealth[scale=0.8]}, black, thick] (0,0) -- (4.5,0) node[right] {$V$};
    \draw[-{Stealth[scale=0.8]}, black, thick] (0,0) -- (0,3.5) node[above] {$p$};
    
    \coordinate (1) at (1,1);
    \coordinate (2) at (3,2);
    \coordinate (3) at (3,1);
    \coordinate (4) at (1,2);
    
    \draw[ultra thick, brandPrimary, ->, >=stealth] (1) -- (2);
    \draw[ultra thick, brandPrimary, ->, >=stealth] (2) -- (3);
    \draw[ultra thick, brandPrimary, ->, >=stealth] (3) -- (1);
    
    \fill[black] (1) circle (1.5pt) node[below left] {4};
    \fill[black] (2) circle (1.5pt) node[above right] {2};
    \fill[black] (3) circle (1.5pt) node[below right] {3};
    \fill[black] (1) circle (1.5pt) ++(0,1) node[above left] {1};
    \draw[ultra thick, brandPrimary] (1,1) -- (1,2) -- (3,2);
    
    \node[left] at (0,1) {$p_0$};
    \node[left] at (0,2) {$2p_0$};
    \node[below] at (1,0) {$V_0$};
    \node[below] at (3,0) {$3V_0$};
    \node[below left] at (0,0) {0};
\end{tikzpicture}
\end{center}

1) Аргон отдаёт холодильнику положительное количество теплоты только в процессе 3–4.\\
2) В процессе 3–4 внутренняя энергия аргона уменьшается.\\
3) Работа аргона за цикл равна $2p_0V_0$.\\
4) Максимальная абсолютная температура аргона в цикле в $6$ раз больше минимальной.\\
5) В процессе 4–1 аргон получает от нагревателя положительное количество теплоты.

\kim{9}
\end{taskbasic}
\answer{24}

% === ЗАДАЧА 4 ===
% Тег: #МКТ #ТЕРМОДИНАМИКА #КИМ_09 #ВАРИАНТ_04
\begin{taskbasic}[code={\label{task:mkt_therm_v4_n9}}]
\textbf{Задача 4.} Один моль разреженного аргона является рабочим телом в тепловом двигателе, который работает по циклу, показанному на рисунке в переменных $p-V$. Из приведённого ниже списка выберите все верные утверждения, характеризующие работу двигателя.

\nopagebreak\smallskip
\begin{center}
\begin{tikzpicture}[scale=1.1, every node/.style={font=\footnotesize}]
    \draw[xstep=1, ystep=1, gray!30, thin] (0,0) grid (4,3);
    \draw[-{Stealth[scale=0.8]}, black, thick] (0,0) -- (4.5,0) node[right] {$V$};
    \draw[-{Stealth[scale=0.8]}, black, thick] (0,0) -- (0,3.5) node[above] {$p$};
    
    \draw[ultra thick, brandPrimary] (1,1) -- (1,2) -- (3,2) -- (3,1) -- cycle;
    \draw[ultra thick, brandPrimary, ->, >=stealth] (1,1) -- (1,1.6);
    \draw[ultra thick, brandPrimary, ->, >=stealth] (1,2) -- (2.2,2);
    \draw[ultra thick, brandPrimary, ->, >=stealth] (3,2) -- (3,1.4);
    \draw[ultra thick, brandPrimary, ->, >=stealth] (3,1) -- (1.8,1);
    
    \fill[black] (1,2) circle (1.5pt) node[above left] {1};
    \fill[black] (2,2) circle (0pt); 
    \fill[black] (3,2) circle (1.5pt) node[above right] {2};
    \fill[black] (3,1) circle (1.5pt) node[below right] {3};
    \fill[black] (1,1) circle (1.5pt) node[below left] {4};
    
    \node[left] at (0,1) {$p_0$};
    \node[left] at (0,2) {$2p_0$};
    \node[below] at (1,0) {$V_0$};
    \node[below] at (3,0) {$3V_0$};
    \node[below left] at (0,0) {0};
\end{tikzpicture}
\end{center}

1) Аргон получает положительное количество теплоты от нагревателя только в процессах 1–2 и 4–1.\\
2) В процессе 3–4 внутренняя энергия аргона не изменяется.\\
3) Работа аргона за цикл равна $6p_0V_0$.\\
4) Максимальная абсолютная температура аргона в цикле в $6$ раз больше минимальной.\\
5) В процессе 4–1 аргон отдаёт холодильнику положительное количество теплоты.

\kim{9}
\end{taskbasic}
\answer{15}
% ====================================================================
% ФАЙЛ: tasks/02_mkt/mkt_thermodynamics_first_law_part2.tex
% ТЕМА: ПЕРВОЕ НАЧАЛО ТЕРМОДИНАМИКИ. РАБОТА ГАЗА. ВНУТРЕННЯЯ ЭНЕРГИЯ
% ПАЧКА 2: ВАРИАНТЫ 6–12
% ====================================================================

% === ЗАДАЧА 1 ===
% Тег: #МКТ #ТЕРМОДИНАМИКА #КИМ_08 #ВАРИАНТ_11
\begin{taskbasic}[code={\label{task:mkt_therm_v11_n8}}]
\textbf{Задача 1.} На рисунке показан график циклического процесса, проведенного с идеальным одноатомным газом, в координатах $p-V$. В ходе изобарного расширения 1–2 газ совершил работу $40$~Дж. Какую работу совершили внешние силы над газом в процессе 3–1?

\nopagebreak\smallskip
\begin{center}
\begin{tikzpicture}[scale=1.1, every node/.style={font=\footnotesize}]
    \draw[xstep=1, ystep=1, gray!30, thin] (0,0) grid (3,3);
    \draw[-{Stealth[scale=0.8]}, black, thick] (0,0) -- (3.5,0) node[right] {$V$};
    \draw[-{Stealth[scale=0.8]}, black, thick] (0,0) -- (0,3.5) node[above] {$p$};
    
    \coordinate (1) at (1,2);
    \coordinate (2) at (2,2);
    \coordinate (3) at (2,1);
    
    \draw[ultra thick, brandPrimary, ->, >=stealth] (1) -- (1.6,2);
    \draw[ultra thick, brandPrimary] (1.5,2) -- (2);
    \draw[ultra thick, brandPrimary, ->, >=stealth] (2) -- (2,1.4);
    \draw[ultra thick, brandPrimary] (2,1.5) -- (3);
    \draw[ultra thick, brandPrimary, ->, >=stealth] (3) -- (1.5,1.5);
    \draw[ultra thick, brandPrimary] (1.6,1.6) -- (1);
    
    \draw[dashed, gray] (0,1) -- (2,1);
    \draw[dashed, gray] (0,2) -- (1,2);
    \draw[dashed, gray] (1,0) -- (1,2);
    \draw[dashed, gray] (2,0) -- (2,2);
    
    \fill[black] (1) circle (1.5pt) node[above left] {1};
    \fill[black] (2) circle (1.5pt) node[above right] {2};
    \fill[black] (3) circle (1.5pt) node[below right] {3};
    
    \node[left] at (0,1) {$p_0$};
    \node[left] at (0,2) {$2p_0$};
    \node[below] at (1,0) {$V_0$};
    \node[below] at (2,0) {$2V_0$};
    \node[below left] at (0,0) {0};
\end{tikzpicture}
\end{center}

\kim{8}
\end{taskbasic}
\answer{30}

% === ЗАДАЧА 2 ===
% Тег: #МКТ #ТЕРМОДИНАМИКА #КИМ_08 #ВАРИАНТ_12
\begin{taskbasic}[code={\label{task:mkt_therm_v12_n8}}]
\textbf{Задача 2.} На рисунке показан график циклического процесса, проведенного с идеальным одноатомным газом, в координатах $p-V$. В ходе изобарного расширения 1–2 газ совершил работу $80$~Дж. Какую работу совершили внешние силы над газом в процессе 3–1?

\nopagebreak\smallskip
\begin{center}
\begin{tikzpicture}[scale=1.1, every node/.style={font=\footnotesize}]
    \draw[xstep=1, ystep=1, gray!30, thin] (0,0) grid (3,3);
    \draw[-{Stealth[scale=0.8]}, black, thick] (0,0) -- (3.5,0) node[right] {$V$};
    \draw[-{Stealth[scale=0.8]}, black, thick] (0,0) -- (0,3.5) node[above] {$p$};
    
    \coordinate (1) at (1,2);
    \coordinate (2) at (2,2);
    \coordinate (3) at (2,1);
    
    \draw[ultra thick, brandPrimary, ->, >=stealth] (1) -- (1.6,2);
    \draw[ultra thick, brandPrimary] (1.5,2) -- (2);
    \draw[ultra thick, brandPrimary, ->, >=stealth] (2) -- (2,1.4);
    \draw[ultra thick, brandPrimary] (2,1.5) -- (3);
    \draw[ultra thick, brandPrimary, ->, >=stealth] (3) -- (1.5,1.5);
    \draw[ultra thick, brandPrimary] (1.6,1.6) -- (1);
    
    \draw[dashed, gray] (0,1) -- (2,1);
    \draw[dashed, gray] (0,2) -- (1,2);
    \draw[dashed, gray] (1,0) -- (1,2);
    \draw[dashed, gray] (2,0) -- (2,2);
    
    \fill[black] (1) circle (1.5pt) node[above left] {1};
    \fill[black] (2) circle (1.5pt) node[above right] {2};
    \fill[black] (3) circle (1.5pt) node[below right] {3};
    
    \node[left] at (0,1) {$p_0$};
    \node[left] at (0,2) {$2p_0$};
    \node[below] at (1,0) {$V_0$};
    \node[below] at (2,0) {$2V_0$};
    \node[below left] at (0,0) {0};
\end{tikzpicture}
\end{center}

\kim{8}
\end{taskbasic}
\answer{60}

% ====================================================================
% ФАЙЛ: tasks/02_mkt/mkt_thermodynamics_first_law_part3.tex
% ТЕМА: ПЕРВОЕ НАЧАЛО ТЕРМОДИНАМИКИ. РАБОТА ГАЗА. ВНУТРЕННЯЯ ЭНЕРГИЯ
% ПАЧКА 3: ВАРИАНТЫ 13–22
% ====================================================================

% === ЗАДАЧА 1 ===
% Тег: #МКТ #ТЕРМОДИНАМИКА #КИМ_08 #ВАРИАНТ_13
\begin{taskbasic}[code={\label{task:mkt_therm_v13_n8}}]
\textbf{Задача 1.} На рисунке показано, как меняется давление идеального газа в зависимости от его объёма при переходе из состояния 1 в состояние 2, а затем в состояние 3. В ходе процесса 2–3 газ совершил работу $360$~Дж. Чему равна работа газа в этом процессе?

\nopagebreak\smallskip
\begin{center}
\begin{tikzpicture}[scale=1.0, every node/.style={font=\footnotesize}]
    \draw[xstep=1, ystep=1, gray!30, thin] (0,0) grid (4,4);
    \draw[-{Stealth[scale=0.8]}, black, thick] (0,0) -- (4.5,0) node[right] {$V$};
    \draw[-{Stealth[scale=0.8]}, black, thick] (0,0) -- (0,4.5) node[above] {$p$};
    
    \coordinate (1) at (1,1);
    \coordinate (2) at (1,3);
    \coordinate (3) at (3,1);
    
    \draw[ultra thick, brandPrimary, ->, >=stealth] (1) -- (1,2.2);
    \draw[ultra thick, brandPrimary] (1,2.0) -- (2);
    \draw[ultra thick, brandPrimary, ->, >=stealth] (2) -- (2.2,1.8);
    \draw[ultra thick, brandPrimary] (2.1,1.9) -- (3);
    
    \draw[dashed, gray] (0,1) -- (1,1);
    \draw[dashed, gray] (0,3) -- (1,3);
    \draw[dashed, gray] (1,0) -- (1,1);
    \draw[dashed, gray] (3,0) -- (3,1);
    \draw[dashed, gray] (2,0) -- (2,2);
    
    \fill[black] (1) circle (1.5pt) node[below right] {1};
    \fill[black] (2) circle (1.5pt) node[above right] {2};
    \fill[black] (3) circle (1.5pt) node[above right] {3};
    
    \node[left] at (0,1) {$p_0$};
    \node[left] at (0,3) {$2p_0$};
    \node[below] at (1,0) {$V_0$};
    \node[below] at (2,0) {$2V_0$};
    \node[below] at (3,0) {$3V_0$};
    \node[below left] at (0,0) {0};
\end{tikzpicture}
\end{center}

\kim{8}
\end{taskbasic}
\answer{360}

% === ЗАДАЧА 2 ===
% Тег: #МКТ #ТЕРМОДИНАМИКА #КИМ_08 #ВАРИАНТ_14
\begin{taskbasic}[code={\label{task:mkt_therm_v14_n8}}]
\textbf{Задача 2.} На рисунке показано, как меняется давление идеального газа в зависимости от его объёма при переходе из состояния 1 в состояние 2, а затем в состояние 3. В ходе процесса 2–3 газ совершил работу $360$~Дж. Чему равна величина произведения $p_0V_0$?

\nopagebreak\smallskip
\begin{center}
\begin{tikzpicture}[scale=1.0, every node/.style={font=\footnotesize}]
    \draw[xstep=1, ystep=1, gray!30, thin] (0,0) grid (4,4);
    \draw[-{Stealth[scale=0.8]}, black, thick] (0,0) -- (4.5,0) node[right] {$V$};
    \draw[-{Stealth[scale=0.8]}, black, thick] (0,0) -- (0,4.5) node[above] {$p$};
    
    \coordinate (1) at (1,1);
    \coordinate (2) at (1,3);
    \coordinate (3) at (3,1);
    
    \draw[ultra thick, brandPrimary, ->, >=stealth] (1) -- (1,2.2);
    \draw[ultra thick, brandPrimary] (1,2.0) -- (2);
    \draw[ultra thick, brandPrimary, ->, >=stealth] (2) -- (2.2,1.8);
    \draw[ultra thick, brandPrimary] (2.1,1.9) -- (3);
    
    \draw[dashed, gray] (0,1) -- (1,1);
    \draw[dashed, gray] (0,3) -- (1,3);
    \draw[dashed, gray] (1,0) -- (1,1);
    \draw[dashed, gray] (3,0) -- (3,1);
    \draw[dashed, gray] (2,0) -- (2,2);
    
    \fill[black] (1) circle (1.5pt) node[below right] {1};
    \fill[black] (2) circle (1.5pt) node[above right] {2};
    \fill[black] (3) circle (1.5pt) node[above right] {3};
    
    \node[left] at (0,1) {$p_0$};
    \node[left] at (0,3) {$2p_0$};
    \node[below] at (1,0) {$V_0$};
    \node[below] at (2,0) {$2V_0$};
    \node[below] at (3,0) {$3V_0$};
    \node[below left] at (0,0) {0};
\end{tikzpicture}
\end{center}

\kim{8}
\end{taskbasic}
\answer{120}

% === ЗАДАЧА 3 ===
% Тег: #МКТ #ТЕРМОДИНАМИКА #КИМ_08 #ВАРИАНТ_19
\begin{taskbasic}[code={\label{task:mkt_therm_v19_n8}}]
\textbf{Задача 3.} На рисунке показано, как меняется давление идеального газа в зависимости от его объёма при переходе из состояния 1 в состояние 2, а затем в состояние 3. В ходе процесса 2–3 газ совершил работу $1000$~Дж. Чему равна работа газа в процессе 1–2?

\nopagebreak\smallskip
\begin{center}
\begin{tikzpicture}[scale=1.0, every node/.style={font=\footnotesize}]
    \draw[xstep=1, ystep=1, gray!30, thin] (0,0) grid (4,4);
    \draw[-{Stealth[scale=0.8]}, black, thick] (0,0) -- (4.5,0) node[right] {$V$};
    \draw[-{Stealth[scale=0.8]}, black, thick] (0,0) -- (0,4.5) node[above] {$p$};
    
    \coordinate (1) at (1,3);
    \coordinate (2) at (2,3);
    \coordinate (3) at (3,1);
    
    \draw[ultra thick, brandPrimary, ->, >=stealth] (1) -- (1.6,3);
    \draw[ultra thick, brandPrimary] (1.5,3) -- (2);
    \draw[ultra thick, brandPrimary, ->, >=stealth] (2) -- (2.5,2);
    \draw[ultra thick, brandPrimary] (2.4,2.2) -- (3);
    
    \draw[dashed, gray] (0,1) -- (3,1);
    \draw[dashed, gray] (0,3) -- (1,3);
    \draw[dashed, gray] (1,0) -- (1,3);
    \draw[dashed, gray] (2,0) -- (2,3);
    \draw[dashed, gray] (3,0) -- (3,1);
    
    \fill[black] (1) circle (1.5pt) node[above left] {1};
    \fill[black] (2) circle (1.5pt) node[above right] {2};
    \fill[black] (3) circle (1.5pt) node[above right] {3};
    
    \node[left] at (0,1) {$p_0$};
    \node[left] at (0,3) {$3p_0$};
    \node[below] at (1,0) {$V_0$};
    \node[below] at (2,0) {$2V_0$};
    \node[below] at (3,0) {$3V_0$};
    \node[below left] at (0,0) {0};
\end{tikzpicture}
\end{center}

\kim{8}
\end{taskbasic}
\answer{1500}

% === ЗАДАЧА 4 ===
% Тег: #МКТ #ТЕРМОДИНАМИКА #КИМ_21 #ВАРИАНТ_21
\begin{taskbasic}[code={\label{task:mkt_therm_v21_n21}}]
\textbf{Задача 4.} Один моль одноатомного идеального газа участвует в циклическом процессе 1–2–3–4–1, график которого изображён на рисунке в координатах $V-T$, где $V$ — объём газа, $T$ — абсолютная температура. Опираясь на законы молекулярной физики и термодинамики, сравните работу газа в процессе 2–3 и работу внешних сил в процессе 4–1. Постройте график цикла в координатах $p-V$, где $p$ — давление газа, $V$ — объём газа.

\nopagebreak\smallskip
\begin{center}
\begin{tikzpicture}[scale=1.0, every node/.style={font=\footnotesize}]
    \draw[xstep=1, ystep=1, gray!30, thin] (0,0) grid (7,4);
    \draw[-{Stealth[scale=0.8]}, black, thick] (0,0) -- (7.5,0) node[right] {$T$};
    \draw[-{Stealth[scale=0.8]}, black, thick] (0,0) -- (0,4.5) node[above] {$V$};
    
    \coordinate (1) at (1,1);
    \coordinate (2) at (6,1);
    \coordinate (3) at (6,3);
    \coordinate (4) at (3,3);
    
    \draw[dashed, thin, gray] (0,0) -- (1);
    \draw[dashed, thin, gray] (0,0) -- (4);
    
    \draw[ultra thick, brandPrimary, ->, >=stealth] (1) -- (4);
    \draw[ultra thick, brandPrimary, ->, >=stealth] (4) -- (3);
    \draw[ultra thick, brandPrimary, ->, >=stealth] (3) -- (2);
    \draw[ultra thick, brandPrimary, ->, >=stealth] (2) -- (1);
    
    \fill[black] (1) circle (1.5pt) node[above left] {1};
    \fill[black] (2) circle (1.5pt) node[below right] {2};
    \fill[black] (3) circle (1.5pt) node[above right] {3};
    \fill[black] (4) circle (1.5pt) node[above left] {4};
    
    \node[left] at (0,1) {$V_0$};
    \node[left] at (0,3) {$3V_0$};
    \node[below] at (1,0) {$T_0$};
    \node[below] at (3,0) {$3T_0$};
    \node[below] at (6,0) {$6T_0$};
    \node[below left] at (0,0) {0};
\end{tikzpicture}
\end{center}

\kim{21}
\end{taskbasic}
\answer{Сравнение работы}

% === ЗАДАЧА 5 ===
% Тег: #МКТ #ТЕРМОДИНАМИКА #КИМ_21 #ВАРИАНТ_22
\begin{taskbasic}[code={\label{task:mkt_therm_v22_n21}}]
\textbf{Задача 5.} Один моль одноатомного идеального газа участвует в циклическом процессе 1–2–3–4–1, график которого изображён на рисунке в координатах $p-T$, где $p$ — давление газа, $T$ — абсолютная температура. Опираясь на законы молекулярной физики и термодинамики, сравните модуль работы газа в процессе 2–3 и модуль работы внешних сил в процессе 4–1. Постройте график цикла в координатах $p-V$, где $p$ — давление газа, $V$ — объём газа.

\nopagebreak\smallskip
\begin{center}
\begin{tikzpicture}[scale=1.0, every node/.style={font=\footnotesize}]
    \draw[xstep=1, ystep=1, gray!30, thin] (0,0) grid (7,6);
    \draw[-{Stealth[scale=0.8]}, black, thick] (0,0) -- (7.5,0) node[right] {$T$};
    \draw[-{Stealth[scale=0.8]}, black, thick] (0,0) -- (0,6.5) node[above] {$p$};
    
    \coordinate (1) at (1,1);
    \coordinate (2) at (6,5);
    \coordinate (3) at (6,3);
    \coordinate (4) at (2,1);
    
    \draw[dashed, thin, gray] (0,0) -- (1);
    \draw[dashed, thin, gray] (0,0) -- (3);
    
    \draw[ultra thick, brandPrimary, ->, >=stealth] (1) -- (2);
    \draw[ultra thick, brandPrimary, ->, >=stealth] (2) -- (3);
    \draw[ultra thick, brandPrimary, ->, >=stealth] (3) -- (4);
    \draw[ultra thick, brandPrimary, ->, >=stealth] (4) -- (1);
    
    \fill[black] (1) circle (1.5pt) node[above left] {1};
    \fill[black] (2) circle (1.5pt) node[above right] {2};
    \fill[black] (3) circle (1.5pt) node[right=2pt] {3};
    \fill[black] (4) circle (1.5pt) node[below right] {4};
    
    \node[left] at (0,1) {$p_0$};
    \node[left] at (0,3) {$3p_0$};
    \node[left] at (0,5) {$5p_0$};
    \node[below] at (1,0) {$T_0$};
    \node[below] at (2,0) {$2T_0$};
    \node[below] at (6,0) {$6T_0$};
    \node[below left] at (0,0) {0};
\end{tikzpicture}
\end{center}

\kim{21}
\end{taskbasic}
\answer{Сравнение работы}

% ====================================================================
% ФАЙЛ: tasks/02_mkt/mkt_thermodynamics_first_law_part4.tex
% ТЕМА: ПЕРВОЕ НАЧАЛО ТЕРМОДИНАМИКИ. РАБОТА ГАЗА. ВНУТРЕННЯЯ ЭНЕРГИЯ
% ПАЧКА 4: ВАРИАНТЫ 23–30
% ====================================================================

% === ЗАДАЧА 1 ===
% Тег: #МКТ #ТЕРМОДИНАМИКА #КИМ_08 #ВАРИАНТ_25
\begin{taskbasic}[code={\label{task:mkt_therm_v25_n8}}]
\textbf{Задача 1.} На рисунке показано, как меняется давление идеального газа в зависимости от его объёма при переходе из состояния 1 в состояние 2, а затем в состояние 3. В ходе процесса 2–3 газ совершил работу $600$~Дж. Чему равна работа газа, если бы он расширялся изобарно при давлении $2p_0$ от объёма $V_0$ до $2V_0$?

\nopagebreak\smallskip
\begin{center}
\begin{tikzpicture}[scale=1.0, every node/.style={font=\footnotesize}]
    \draw[xstep=1, ystep=1, gray!30, thin] (0,0) grid (4,4);
    \draw[-{Stealth[scale=0.8]}, black, thick] (0,0) -- (4.5,0) node[right] {$V$};
    \draw[-{Stealth[scale=0.8]}, black, thick] (0,0) -- (0,4.5) node[above] {$p$};
    
    \coordinate (1) at (1,1);
    \coordinate (2) at (1,3);
    \coordinate (3) at (3,1);
    
    \draw[ultra thick, brandPrimary, ->, >=stealth] (1) -- (1,2.2);
    \draw[ultra thick, brandPrimary] (1,2.0) -- (2);
    \draw[ultra thick, brandPrimary, ->, >=stealth] (2) -- (2.2,1.8);
    \draw[ultra thick, brandPrimary] (2.1,1.9) -- (3);
    
    \draw[dashed, gray] (0,1) -- (1,1);
    \draw[dashed, gray] (0,3) -- (1,3);
    \draw[dashed, gray] (1,0) -- (1,1);
    \draw[dashed, gray] (3,0) -- (3,1);
    
    \fill[black] (1) circle (1.5pt) node[below right] {1};
    \fill[black] (2) circle (1.5pt) node[above right] {2};
    \fill[black] (3) circle (1.5pt) node[above right] {3};
    
    \node[left] at (0,1) {$p_0$};
    \node[left] at (0,3) {$2p_0$};
    \node[below] at (1,0) {$V_0$};
    \node[below] at (2,0) {$2V_0$};
    \node[below] at (3,0) {$3V_0$};
    \node[below left] at (0,0) {0};
\end{tikzpicture}
\end{center}

\kim{8}
\end{taskbasic}
\answer{400}

% === ЗАДАЧА 2 ===
% Тег: #МКТ #ТЕРМОДИНАМИКА #КИМ_08 #ВАРИАНТ_26
\begin{taskbasic}[code={\label{task:mkt_therm_v26_n8}}]
\textbf{Задача 2.} На рисунке показано, как меняется давление идеального газа в зависимости от его объёма при переходе из состояния 1 in состояние 2, а затем в состояние 3. В ходе процесса 2–3 газ совершил работу $600$~Дж. Чему равна величина произведения $p_0V_0$?

\nopagebreak\smallskip
\begin{center}
\begin{tikzpicture}[scale=1.0, every node/.style={font=\footnotesize}]
    \draw[xstep=1, ystep=1, gray!30, thin] (0,0) grid (4,4);
    \draw[-{Stealth[scale=0.8]}, black, thick] (0,0) -- (4.5,0) node[right] {$V$};
    \draw[-{Stealth[scale=0.8]}, black, thick] (0,0) -- (0,4.5) node[above] {$p$};
    
    \coordinate (1) at (1,1);
    \coordinate (2) at (1,3);
    \coordinate (3) at (3,1);
    
    \draw[ultra thick, brandPrimary, ->, >=stealth] (1) -- (1,2.2);
    \draw[ultra thick, brandPrimary] (1,2.0) -- (2);
    \draw[ultra thick, brandPrimary, ->, >=stealth] (2) -- (2.2,1.8);
    \draw[ultra thick, brandPrimary] (2.1,1.9) -- (3);
    
    \draw[dashed, gray] (0,1) -- (1,1);
    \draw[dashed, gray] (0,3) -- (1,3);
    \draw[dashed, gray] (1,0) -- (1,1);
    \draw[dashed, gray] (3,0) -- (3,1);
    
    \fill[black] (1) circle (1.5pt) node[below right] {1};
    \fill[black] (2) circle (1.5pt) node[above right] {2};
    \fill[black] (3) circle (1.5pt) node[above right] {3};
    
    \node[left] at (0,1) {$p_0$};
    \node[left] at (0,3) {$2p_0$};
    \node[below] at (1,0) {$V_0$};
    \node[below] at (2,0) {$2V_0$};
    \node[below] at (3,0) {$3V_0$};
    \node[below left] at (0,0) {0};
\end{tikzpicture}
\end{center}

\kim{8}
\end{taskbasic}
\answer{200}